\documentclass[11pt,a4paper]{article}

% ------------------------------------------------------------------
% Paquetes (todos estándar, compatibles con Overleaf sin configuración extra)
% ------------------------------------------------------------------
\usepackage[utf8]{inputenc}
\usepackage[T1]{fontenc}
\usepackage[spanish]{babel}
\usepackage[margin=2.5cm]{geometry}
\usepackage{xcolor}
\usepackage{booktabs}
\usepackage{longtable}
\usepackage{array}
\usepackage{ragged2e}
\usepackage{enumitem}
\usepackage{hyperref}
\usepackage{fancyhdr}
\usepackage{titlesec}
\usepackage{colortbl}
\usepackage{multicol}

\definecolor{critico}{RGB}{178,34,34}
\definecolor{advertencia}{RGB}{184,115,10}
\definecolor{ok}{RGB}{30,120,90}
\definecolor{gris}{RGB}{90,90,90}
\definecolor{fondoheader}{RGB}{240,240,240}

\hypersetup{
  colorlinks=true,
  linkcolor=teal,
  citecolor=teal,
  urlcolor=teal,
  pdftitle={Gamma Lab - Informe de rendimiento (línea base v1.0)},
  pdfauthor={Auditoría de rendimiento}
}

\newcommand{\okmark}{\textcolor{ok}{\textbf{RÁPIDO}}}
\newcommand{\warnmark}{\textcolor{advertencia}{\textbf{A VIGILAR}}}
\newcommand{\critmark}{\textcolor{critico}{\textbf{CUELLO DE BOTELLA}}}
\newcommand{\estmark}{\textcolor{gris}{\textit{[estimado]}}}

\newcolumntype{L}[1]{>{\small\RaggedRight\arraybackslash}p{#1}}

\pagestyle{fancy}
\fancyhf{}
\fancyhead[L]{Gamma Lab -- Informe de rendimiento (línea base v1.0)}
\fancyhead[R]{\thepage}
\renewcommand{\headrulewidth}{0.4pt}

\titleformat{\section}{\normalfont\Large\bfseries}{\thesection}{1em}{}
\titleformat{\subsection}{\normalfont\large\bfseries}{\thesubsection}{1em}{}

\title{\textbf{Gamma Lab}\\[4pt]\large Informe de rendimiento: línea base de Gamma Lab 1.0\\
\large Kernel, plugins, cómputo numérico y renderizado VTK -- pensado para comparar contra Gamma Lab 2.0}
\author{Auditoría de rendimiento}
\date{16 de agosto de 2026}

\begin{document}

\maketitle
\thispagestyle{empty}

\begin{abstract}
\noindent
Este informe documenta el rendimiento medido de \textbf{Gamma Lab 1.0} completo: arranque del kernel y del
sistema de plugins, carga de archivos, corte de trials, preprocesamiento (filtro), cómputo numérico de cada
plugin de análisis (FFT, PSD, Wavelet) y \textbf{todo} el renderizado VTK de la aplicación (los 19 archivos
del repositorio que usan VTK). Ningún archivo de \texttt{core/} o \texttt{plugins/} fue modificado para
producir este informe -- todas las cifras provienen de ejecuciones reales de \texttt{pytest} contra el
código real de la aplicación, con el archivo ABF real de \texttt{test/data/17308005.abf}, sin mockear
ninguna fórmula matemática ni ninguna llamada a VTK. La suite que genera estos números vive en
\texttt{test/performance\_\allowbreak test/} y está explícitamente diseñada para volver a correrse, sin
cambios de metodología, contra \textbf{Gamma Lab 2.0} durante su desarrollo este mes -- la
sección~\ref{sec:comparacion-v2} explica el procedimiento exacto.
\end{abstract}

\tableofcontents
\newpage

% ==================================================================
\section{Resumen ejecutivo}
% ==================================================================

Gamma Lab 1.0, medido en la máquina de referencia de la sección~\ref{sec:entorno}, tarda:

\begin{itemize}[nosep]
  \item \textbf{$\approx$20\,ms} en arrancar el kernel y cargar los 16 plugins (irrelevante frente al resto).
  \item \textbf{$\approx$78\,ms} en cargar un archivo ABF de 1.8\,M muestras y cortarlo en 60 trials.
  \item \textbf{$\approx$1.02\,s} en total desde que el usuario pide un escalograma de wavelet de un solo
        trial hasta que aparece en pantalla -- de los cuales el \textbf{95\%} son solo dos pasos: el CWT
        (216\,ms) y el render del mapa 2D (725\,ms).
  \item \textbf{$\approx$12.6\,s} (estimado por extrapolación lineal) en calcular el escalograma promedio de
        una sesión típica de 60 trials en \texttt{wavelet\_\allowbreak average} -- ya corre en un
        \texttt{QThread} y no congela la interfaz, pero sigue siendo el número más grande de toda la
        aplicación.
\end{itemize}

\textbf{El hallazgo central es el mismo patrón, repetido en al menos cinco lugares del código:} construir un
arreglo VTK con un bucle de Python (\texttt{InsertNextValue}/\texttt{SetScalarComponentFromFloat} punto por
punto) en vez de vectorizar con \texttt{numpy\_\allowbreak to\_\allowbreak vtk}, algo que \textbf{ya existe,
ya funciona y ya se usa en 9 de los 19 archivos que usan VTK en el repositorio}
(\texttt{core/utils/adapters.py} y sus consumidores). El micro-benchmark sintético de la
sección~\ref{sec:micro-bench} mide una diferencia de \textbf{65x a 150x} entre ambos patrones para el mismo
volumen de datos. Vectorizar los cuatro puntos que faltan (\texttt{wavelet}, \texttt{wavelet\_\allowbreak
average}, \texttt{erp} para su mapa 2D, y \texttt{filter} para su gráfico de línea) es, con estos números
sobre la mesa, la optimización de mayor impacto por esfuerzo disponible para Gamma Lab 2.0.

\begin{center}
\begin{tabular}{lrrl}
\toprule
\textbf{Cuello de botella} & \textbf{Costo medido} & \textbf{Hilo} & \textbf{Prioridad} \\
\midrule
\texttt{wavelet\_\allowbreak average} -- CWT $\times$60 trials & $\approx$12.6\,s \estmark & QThread & \warnmark \\
\texttt{wavelet}/\texttt{wavelet\_\allowbreak average} -- render 2D & 725--742\,ms & Principal & \critmark \\
\texttt{wavelet} -- CWT (1 trial) & 216\,ms & Principal & \critmark \\
\texttt{filter} -- render de línea (señal completa) & $\approx$336\,ms \estmark & Principal & \critmark \\
\texttt{fft\_\allowbreak average}/\texttt{psd\_\allowbreak average} -- cómputo & 9--11\,ms & Principal & \warnmark \\
Arranque del kernel y los 16 plugins & 20\,ms & Principal & \okmark \\
\bottomrule
\end{tabular}
\end{center}

% ==================================================================
\section{Objetivo y alcance}
% ==================================================================

Este informe tiene un único propósito: dejar una \textbf{línea base numérica y reproducible} del rendimiento
de Gamma Lab 1.0 completo, para que dentro de este mismo mes -- mientras se desarrolla Gamma Lab 2.0 -- se
pueda responder con datos, no con impresiones, a la pregunta ``¿la versión nueva es más rápida?''.

Por eso el alcance es deliberadamente amplio: no se midió solo el renderizado VTK (que ya se había medido
parcialmente en \texttt{docs/reporte\_\allowbreak tests.tex}, sección de rendimiento), sino \textbf{el kernel},
\textbf{el sistema de plugins}, \textbf{la carga de archivos}, \textbf{el corte de trials}, \textbf{el
preprocesamiento} y \textbf{el cómputo de cada plugin de análisis}, además de repasar y ampliar la parte de
VTK con un inventario completo de los 19 archivos del repositorio que lo usan (sección~\ref{sec:inventario-vtk}).

\subsection{Qué NO cubre este informe}

\begin{itemize}[nosep]
  \item \textbf{Corrección numérica} de los plugins (si el resultado matemático es el correcto) -- eso ya
        está cubierto en \texttt{docs/reporte\_\allowbreak tests.tex}, que valida cada plugin contra una
        referencia de MATLAB. Este informe asume que el código hace lo que dice hacer y mide únicamente
        cuanto tarda en hacerlo.
  \item La construcción de \texttt{MainWindow} (la ventana principal) -- requiere una \texttt{QApplication}
        con un ciclo de eventos completo, no solo un stack VTK offscreen; ver
        sección~\ref{sec:limitaciones}.
  \item \texttt{FileIOService.load\_\allowbreak edf}/\texttt{load\_\allowbreak mat} -- no hay ningún archivo
        \texttt{.edf}/\texttt{.mat} de referencia en \texttt{test/data/}.
\end{itemize}

% ==================================================================
\section{Metodología}
% ==================================================================

La suite completa vive en \texttt{test/performance\_\allowbreak test/} y sigue exactamente la misma filosofía
que el resto de la suite de pruebas del proyecto: \textbf{código real, datos reales, sin mocks}.

\begin{enumerate}
  \item \textbf{Datos reales:} todos los benchmarks que necesitan una señal usan el mismo archivo ABF real,
        \texttt{test/data/17308005.abf} (2 canales, 1\,800\,000 muestras cada uno, 600\,kHz$\to$600\,s a
        3\,kHz de tasa de muestreo), cortado en 60 trials con los mismos parámetros que
        \texttt{test/plugins\_\allowbreak test/} (canal 0, canal de estímulo 1, umbral 0.7, ventana
        $t_0=-0.05\,s$ a $t_1=4.00\,s$, modo \texttt{until\_\allowbreak next\_\allowbreak onset}).
  \item \textbf{Código real, sin mockear:} cada benchmark llama directamente al método real del plugin o
        servicio bajo prueba (\texttt{plug.\_\allowbreak compute\_\allowbreak fft}, \texttt{plug.compute\_\allowbreak
        wavelet}, \texttt{plug.render\_\allowbreak scalogram}, \texttt{FileIOService.load\_\allowbreak abf},
        \texttt{core.plugins.loader.discover}, etc.) -- ninguna fórmula matemática ni ninguna llamada a VTK
        está sustituida por un doble (\textit{mock}).
  \item \textbf{Cronometraje:} \texttt{bench\_\allowbreak common.time\_\allowbreak block(fn, reps=N,
        warmup=W)} ejecuta la función \texttt{W} veces de calentamiento (se descartan, para no medir el
        costo de la primera importación/JIT/cache) y luego \texttt{N} veces cronometradas con
        \texttt{time.perf\_\allowbreak counter()} (reloj monotónico de alta resolución). Se reporta mínimo,
        mediana, media y máximo; este informe usa siempre la \textbf{mediana} como cifra representativa, por
        ser más robusta a un outlier puntual que la media.
  \item \textbf{Persistencia con etiqueta:} cada corrida completa de la suite se guarda en
        \texttt{results/perf\_\allowbreak results.json} bajo una etiqueta (\texttt{GAMMA\_\allowbreak PERF\_\allowbreak
        LABEL}). Los números de este informe corresponden a la etiqueta
        \texttt{gammalab1\_\allowbreak baseline\_\allowbreak 2026\allowbreak-08\allowbreak-16}, generada en la
        misma sesión en la que se escribió este documento.
  \item \textbf{Snapshot de entorno:} cada corrida guarda automáticamente que SO, CPU, Python y versiones de
        librerías (VTK, numpy, scipy, Qt) se usaron (\texttt{bench\_\allowbreak common.record\_\allowbreak
        environment}), justamente para poder juzgar después si una diferencia de tiempos entre v1 y v2 es un
        cambio de código o un cambio de máquina.
\end{enumerate}

% ==================================================================
\section{Entorno de medición}
\label{sec:entorno}
% ==================================================================

Todas las cifras de este informe se midieron en una única corrida, en la siguiente máquina y con las
siguientes versiones de librerías (capturado automáticamente por
\texttt{bench\_\allowbreak common.system\_\allowbreak info()} y por \texttt{Get-\allowbreak CimInstance} de
PowerShell):

\begin{center}
\begin{tabular}{ll}
\toprule
\textbf{Componente} & \textbf{Valor} \\
\midrule
CPU & 13th Gen Intel(R) Core(TM) i7-13620H \\
Núcleos físicos / lógicos & 10 / 16 \\
RAM & 31.7\,GB \\
Sistema operativo & Windows 11 Home, build 26200 (64 bits) \\
Python & 3.11.9 (entorno virtual del proyecto, \texttt{.venv}) \\
VTK & 9.5.2 \\
PyQt5 (Qt) & 5.15.2 \\
numpy & 2.3.4 \\
scipy & 1.16.2 \\
PyWavelets & 1.8.0 \\
pyabf & 2.3.8 \\
\bottomrule
\end{tabular}
\end{center}

\textbf{Por qué importa dejar esto por escrito:} un número de rendimiento sin su entorno es casi inútil para
comparar. Si Gamma Lab 2.0 se mide en otra máquina, con otra versión de VTK o de numpy, cualquier diferencia
observada puede deberse al cambio de entorno y no al cambio de código. La
sección~\ref{sec:comparacion-v2} retoma este punto.

\newpage
% ==================================================================
\section{Inventario completo de VTK en Gamma Lab}
\label{sec:inventario-vtk}
% ==================================================================

Como este informe busca cubrir ``todo el VTK de Gamma Lab'', esta sección enumera los \textbf{19 archivos}
del repositorio que importan \texttt{vtk}/\texttt{vtkmodules}, clasificados según el patrón que usan para
convertir datos de NumPy a estructuras VTK.

\subsection{Patrón vectorizado (\texttt{numpy\_\allowbreak to\_\allowbreak vtk}) -- 10 archivos}

Construyen el arreglo VTK de una sola vez con \texttt{vtk.util.numpy\_\allowbreak support.numpy\_\allowbreak
to\_\allowbreak vtk}, sin iterar en Python. Es el patrón rápido, medido en la
sección~\ref{sec:resultados-render-vectorizado}.

\begin{multicols}{2}
\begin{itemize}[nosep]
  \item \texttt{core/utils/adapters.py} (el helper compartido: \texttt{dataset\_\allowbreak to\_\allowbreak
        vtk\_\allowbreak table}, \texttt{trials\_\allowbreak matrix\_\allowbreak to\_\allowbreak vtk\_\allowbreak table})
  \item \texttt{plugins/analysis/frequency/fft/}
  \item \texttt{plugins/analysis/frequency/fft\_\allowbreak average/}
  \item \texttt{plugins/analysis/frequency/psd/}
  \item \texttt{plugins/analysis/frequency/psd\_\allowbreak average/}
  \item \texttt{plugins/analysis/frequency/relative\_\allowbreak psd/}
  \item \texttt{plugins/analysis/time/erp/} (solo su gráfico \textit{butterfly}; el mapa 2D es aparte, ver abajo)
  \item \texttt{plugins/preprocessing/trials/}
  \item \texttt{plugins/io/open\_\allowbreak signal/}
  \item \texttt{plugins/preprocessing/prepare/artifact\_\allowbreak remove/} (usa \texttt{numpy\_\allowbreak
        to\_\allowbreak vtk} directo, sin pasar por \texttt{adapters.py}, pero el mismo patrón vectorizado)
\end{itemize}
\end{multicols}

\subsection{Patrón de bucle Python (\texttt{InsertNextValue} / \texttt{SetScalarComponentFromFloat}) -- 5 archivos}

Construyen el arreglo VTK punto por punto, dentro de un \texttt{for} de Python. Es el patrón lento, medido
en la sección~\ref{sec:resultados-render-lento} y estimado en la sección~\ref{sec:estimaciones}.

\begin{center}
\begin{longtable}{L{5.7cm}L{2.2cm}L{6.6cm}}
\toprule
\textbf{Archivo} & \textbf{Línea} & \textbf{Qué construye / estado} \\
\midrule
\endhead
\texttt{wavelet\_\allowbreak plugin.py} & 334--336 & Mapa 2D del escalograma. \textbf{Medido}: 725\,ms para 3.05\,M puntos. \\
\texttt{wavelet\_\allowbreak average\_\allowbreak plugin.py} & 444--446 & Mapa 2D del escalograma promedio. \textbf{Medido}: 742\,ms para 3.05\,M puntos. \\
\texttt{erp\_\allowbreak plugin.py} & 340--342 & Mapa 2D \textit{butterfly}. \textbf{Medido}: 36\,ms -- protegido por un tope interno (\texttt{MAX\_\allowbreak SAMPLES=2000}) que recorta el eje de tiempo antes del bucle. \\
\texttt{filter\_\allowbreak plugin.py} & 251--254 & Gráfico de línea de la señal filtrada. \textbf{Sin} tope de muestras -- procesa la señal completa. \textbf{No medido en vivo} (requiere un \texttt{QVTKRenderWindowInteractor} real); \textbf{estimado} en $\approx$336\,ms para 1.8\,M muestras (\S\ref{sec:estimaciones}). \\
\texttt{average\_\allowbreak plugin.py} & 110--112 & Gráfico de línea del promedio entre trials. Protegido por el mismo tope \texttt{MAX\_\allowbreak SAMPLES=2000} que \texttt{erp} -- impacto estimado $<$0.5\,ms, no es prioridad. \\
\bottomrule
\end{longtable}
\end{center}

\subsection{Uso de VTK no relacionado con volumen de datos -- 3 archivos}

Usan VTK pero no construyen arreglos de $N$ puntos en un bucle, así que no aplica esta clasificación
rápido/lento: \texttt{core/utils/vtk\_\allowbreak context\_\allowbreak menu.py} (menu contextual, interacción
de usuario), \texttt{core/services/measurement\_\allowbreak service.py} (mediciones interactivas click A/B) y
\texttt{core/services/export\_\allowbreak service.py} (exportar a imagen/tabla, un archivo por click). Ninguno
de los tres aparece en este informe como benchmark porque su costo está dominado por I/O de disco o por
interacción humana, no por cómputo.

\newpage
% ==================================================================
\section{Resultados: kernel y sistema de plugins}
\label{sec:resultados-kernel}
% ==================================================================

Mide \texttt{core/kernel.py} y \texttt{core/plugins/} reproduciendo exactamente la secuencia de arranque
real de \texttt{main.py} (pasos 1--3, sin la ventana principal -- ver
sección~\ref{sec:limitaciones}): \texttt{PluginManager(plugins\_\allowbreak dir).load\_\allowbreak all()}
(descubre + instancia los 16 plugins reales bajo \texttt{plugins/}) seguido de \texttt{kernel.register\_\allowbreak
plugin(meta.name, plugin)} para cada uno.

\begin{center}
\begin{tabular}{lrr}
\toprule
\textbf{Benchmark} & \textbf{Mediana} & \textbf{Detalle} \\
\midrule
\texttt{plugins.discover} & 19.6\,ms & Escanea \texttt{plugins/}, valida \texttt{properties.yml}, importa cada clase. \\
\texttt{PluginManager.load\_\allowbreak all} & 20.2\,ms & \texttt{discover} + instanciar \texttt{PluginCls(meta)} de los 16 plugins. \\
Arranque completo del kernel & 19.9\,ms & \texttt{load\_\allowbreak all} + \texttt{kernel.register\_\allowbreak plugin} $\times$16 (incluye \texttt{initialize()} de cada plugin). \\
\texttt{Kernel.register\_\allowbreak plugin} (aislado) & 1.4\,ms & Costo propio del kernel con 200 plugins \textit{ficticios} -- 0.007\,ms/plugin. \\
\bottomrule
\end{tabular}
\end{center}

\textbf{Interpretación:} el arranque del kernel \textbf{no es un cuello de botella} -- 20\,ms es
imperceptible comparado con cualquier otra operación de este informe (100$\times$ más rápido que un solo CWT
de wavelet). Además, el propio \texttt{Kernel.register\_\allowbreak plugin} aporta una fracción insignificante
de ese costo (0.007\,ms por plugin, medido de forma aislada con dobles \textit{ficticios} que no importan
PyQt5/VTK/scipy): casi todo el tiempo de arranque es el costo de (re)importar cada módulo de plugin real, que
a su vez está dominado por la primera importación de PyQt5/VTK/scipy en el proceso (compartida entre los 16
plugins, no repetida). \textbf{Recomendación para 2.0:} no invertir esfuerzo de optimización aquí; el
presupuesto de rendimiento rinde mucho más en las secciones siguientes.

% ==================================================================
\section{Resultados: carga de archivos y corte de trials}
% ==================================================================

\begin{center}
\begin{tabular}{lrl}
\toprule
\textbf{Benchmark} & \textbf{Mediana} & \textbf{Tamaño} \\
\midrule
\texttt{FileIOService.load\_\allowbreak abf} & 59.8\,ms & 2 canales $\times$ 1\,800\,000 muestras \\
\texttt{trials.cut\_\allowbreak trials\_\allowbreak single\_\allowbreak channel} & 18.6\,ms & salida 30\,501 $\times$ 60 \\
\bottomrule
\end{tabular}
\end{center}

Juntas, cargar un archivo y generar sus trials toma $\approx$78\,ms -- el usuario lo percibe como
prácticamente instantáneo. No se encontró ninguna operación cuadrática ni ningún recorrido redundante en
ninguno de los dos caminos.

\texttt{FileIOService.load\_\allowbreak edf} y \texttt{load\_\allowbreak mat} no se midieron: \texttt{test/data/}
no tiene ningún archivo \texttt{.edf}/\texttt{.mat} de referencia. Agregar uno (y su benchmark
correspondiente en \texttt{test\_\allowbreak io\_\allowbreak bench.py}) es directo y queda como trabajo
futuro.

% ==================================================================
\section{Resultados: preprocesamiento (filtro)}
% ==================================================================

\texttt{Filter\_\allowbreak plugin.run\_\allowbreak filter} (Butterworth pasa-banda, \texttt{scipy.signal.butter}
+ \texttt{sosfiltfilt}, orden 8, 0.5--4\,Hz) es una función pura, benchmarkeada directamente:

\begin{center}
\begin{tabular}{lrl}
\toprule
\textbf{Benchmark} & \textbf{Mediana} & \textbf{Tamaño} \\
\midrule
\texttt{run\_\allowbreak filter} (señal completa) & 43.2\,ms & 1\,800\,000 muestras \\
\texttt{run\_\allowbreak filter} (1 trial) & 1.8\,ms & 30\,501 muestras \\
\bottomrule
\end{tabular}
\end{center}

El cómputo del filtro en sí es rápido y escala linealmente con el tamaño, como se espera de
\texttt{sosfiltfilt}. El costo real de "Aplicar Filtro" está, en cambio, dominado por su \textbf{renderizado}
-- ver sección~\ref{sec:estimaciones}.

\newpage
% ==================================================================
\section{Resultados: cómputo de los plugins de análisis}
% ==================================================================

Reproduce, con una corrida nueva en la misma máquina que el resto de este informe, la sección de rendimiento
de \texttt{docs/reporte\_\allowbreak tests.tex} (los números de esta sección reemplazan a los de esa
auditoría anterior como referencia vigente). Salvo \texttt{wavelet\_\allowbreak average}, todos
corren hoy en el \textbf{hilo principal de Qt} sin \texttt{QThread} -- la mediana es una medida directa de
cuanto se congela la interfaz por cada clic en ``Calcular''.

\begin{center}
\begin{tabular}{lrl}
\toprule
\textbf{Benchmark} & \textbf{Mediana} & \textbf{Hilo hoy} \\
\midrule
\texttt{fft.\_\allowbreak compute\_\allowbreak fft} (60 trials) & 3.8\,ms & Principal \\
\texttt{fft\_\allowbreak average.\_\allowbreak compute\_\allowbreak fft\_\allowbreak average} & 11.5\,ms & Principal \\
\texttt{psd.\_\allowbreak compute\_\allowbreak psd} (Welch) & 5.3\,ms & Principal \\
\texttt{psd\_\allowbreak average.\_\allowbreak compute\_\allowbreak psd} (Welch) & 9.2\,ms & Principal \\
\texttt{relative\_\allowbreak psd.\_\allowbreak compute\_\allowbreak psd} (Welch) & 6.3\,ms & Principal \\
\texttt{relative\_\allowbreak psd.\_\allowbreak compute\_\allowbreak relative\_\allowbreak psd} (razón) & 0.07\,ms & Principal \\
\texttt{wavelet.compute\_\allowbreak wavelet} (1 trial, CWT) & \textbf{215.6\,ms} & Principal \\
\texttt{wavelet\_\allowbreak average.compute\_\allowbreak wavelet} ($\times$20 trials) & \textbf{4190.2\,ms} & \textit{QThread ya} \\
\bottomrule
\end{tabular}
\end{center}

\textbf{Lectura:} FFT, PSD y PSD promedio son todos submilisegundo a un dígito de milisegundos -- ninguno
justifica moverse a un \texttt{QThread} por sí solo. \texttt{wavelet.compute\_\allowbreak wavelet}, en
cambio, es \textbf{60 veces más lento que el siguiente peor} (\texttt{fft\_\allowbreak average}) y corre en
el hilo principal: cada clic en ``Calcular Wavelet'' congela la interfaz $\approx$216\,ms de forma
perceptible. \texttt{wavelet\_\allowbreak average} escala lo mismo por trial ($\approx$209.5\,ms/trial,
consistente con el número de \texttt{wavelet} individual) pero ya corre en background -- es el número más
grande de todo el informe en términos absolutos ($\approx$4.2\,s para 20 trials, $\approx$12.6\,s
extrapolado a 60), pero al no bloquear la UI su prioridad de arreglo es menor que la de
\texttt{wavelet.compute\_\allowbreak wavelet}, que sí la bloquea.

% ==================================================================
\section{Resultados: renderizado VTK -- patrón lento (bucle Python)}
\label{sec:resultados-render-lento}
% ==================================================================

\begin{center}
\begin{tabular}{lrr}
\toprule
\textbf{Benchmark} & \textbf{Tamaño} & \textbf{Mediana} \\
\midrule
\texttt{wavelet.render\_\allowbreak scalogram} & 998$\times$3051 (3.05\,M puntos) & \textbf{725.3\,ms} \\
\texttt{wavelet\_\allowbreak average.render\_\allowbreak scalogram} & 998$\times$3051 (3.05\,M puntos) & \textbf{741.6\,ms} \\
\texttt{erp.\_\allowbreak render\_\allowbreak heatmap} & 60$\times$30501 $\to$ recortado a 60$\times$1907 & 35.8\,ms \\
\bottomrule
\end{tabular}
\end{center}

Los tres construyen su \texttt{vtkImageData} con \texttt{img.SetScalarComponentFromFloat(i, j, 0, 0, valor)}
dentro de un doble \texttt{for} de Python. \texttt{erp} sale barato en la comparación no porque su código sea
distinto, sino porque un límite interno (\texttt{MAX\_\allowbreak SAMPLES=2000}) ya recorta el eje de tiempo
\emph{antes} de entrar al bucle -- en la práctica itera sobre $\approx$114\,000 puntos, no 1.8\,millones.
\texttt{wavelet} y \texttt{wavelet\_\allowbreak average} no tienen ese límite y pagan el costo completo de 3.05\,M
puntos.

\textbf{Nota de variabilidad medida:} en una corrida anterior de esta misma suite (14 de agosto de 2026, ver
\texttt{results/perf\_\allowbreak results.json}, etiqueta \texttt{v1\_\allowbreak baseline}) este mismo
benchmark de \texttt{wavelet.render\_\allowbreak scalogram} midió 5.70\,s -- casi 8$\times$ más lento que en
la corrida de este informe (725\,ms), sin ningún cambio de código entre ambas. Es una variación atribuible al
primer render offscreen de VTK en un proceso nuevo (inicialización de contexto gráfico), no al código. Al
comparar Gamma Lab 1.0 contra 2.0 (\S\ref{sec:comparacion-v2}) conviene correr cada lado 2--3 veces y mirar
consistencia, no una única corrida.

% ==================================================================
\section{Resultados: renderizado VTK -- patrón vectorizado}
\label{sec:resultados-render-vectorizado}
% ==================================================================

\subsection{Los dos caminos ya vectorizados de \texttt{core/utils/adapters.py}}

\begin{center}
\begin{tabular}{lrr}
\toprule
\textbf{Benchmark} & \textbf{Tamaño} & \textbf{Mediana} \\
\midrule
\texttt{adapters.dataset\_\allowbreak to\_\allowbreak vtk\_\allowbreak table} & 1\,800\,000 filas $\times$ 3 col. & \textbf{21.0\,ms} \\
\texttt{adapters.trials\_\allowbreak matrix\_\allowbreak to\_\allowbreak vtk\_\allowbreak table} & 30\,501 filas $\times$ 61 col. & \textbf{11.1\,ms} \\
\bottomrule
\end{tabular}
\end{center}

Estos dos caminos alimentan los gráficos de línea de 9 de los 16 plugins (sección~\ref{sec:inventario-vtk}) y
son, en volumen de datos equivalente o mayor, \textbf{órdenes de magnitud más rápidos} que el patrón de
bucle: convertir la señal completa de 1.8\,M muestras a una tabla VTK (21\,ms) es más rápido que renderizar un
solo escalograma de 3.05\,M puntos con el patrón de bucle (725\,ms) -- pese a que el volumen de datos es
comparable.

\subsection{Micro-benchmark sintético: bucle vs. vectorizado}
\label{sec:micro-bench}

Para aislar el costo del patrón en sí (sin CWT, sin ABF, solo la conversión a VTK), se construyó el mismo
arreglo de $N$ números aleatorios con \texttt{vtkFloatArray.InsertNextValue} en un bucle de Python, y con
\texttt{numpy\_\allowbreak to\_\allowbreak vtk}:

\begin{center}
\begin{tabular}{rrrr}
\toprule
\textbf{N puntos} & \textbf{Bucle Python} & \textbf{Vectorizado} & \textbf{Factor} \\
\midrule
50\,000  & 4.58\,ms  & 0.031\,ms & \textbf{147.7x} \\
300\,000 & 28.45\,ms & 0.432\,ms & \textbf{65.8x} \\
\bottomrule
\end{tabular}
\end{center}

El factor no es una constante exacta (el camino vectorizado tiene un costo fijo pequeño que pesa más a $N$
chico, y a submilisegundo el ruido de medición es mayor), pero el orden de magnitud es consistente:
\textbf{vectorizar cuesta entre 65 y 150 veces menos} que el bucle de Python, para el mismo volumen de datos,
en esta máquina. Esta cifra es la que sostiene la estimación de la sección siguiente y la recomendación
principal de este informe.

\newpage
% ==================================================================
\section{Estimaciones derivadas (no medidas en vivo)}
\label{sec:estimaciones}
% ==================================================================

Esta sección combina mediciones reales de otras secciones para estimar el costo de dos operaciones que no se
pudieron benchmarkear en vivo por requerir una \texttt{QApplication} completa (ver
sección~\ref{sec:limitaciones}). Se marcan explícitamente \estmark\ y se muestra el cálculo, para que puedan
reemplazarse por una medición real apenas exista un arnés de prueba adecuado.

\subsection{\texttt{filter\_\allowbreak plugin.\_\allowbreak render\_\allowbreak filtered}}

\texttt{filter\_\allowbreak plugin.py:251--254} construye el gráfico de la señal filtrada con un \texttt{for}
que llama dos veces a \texttt{InsertNextValue} por muestra (tiempo y señal), \textbf{sin} ningún tope de
downsampling previo (a diferencia de \texttt{erp}/\texttt{average}, que recortan a 2000 puntos antes de
renderizar). \texttt{on\_apply\_filter} llama a este render \textbf{dos veces} por clic: una para la señal
completa y otra para el trial.

Tomando el costo medido por punto del micro-benchmark
(\S\ref{sec:micro-bench}: $\approx$93\,ns por llamada a \texttt{InsertNextValue} sobre un arreglo, promediado
entre los dos tamaños medidos) y contando 2 llamadas por iteración:

\begin{center}
\begin{tabular}{lrr}
\toprule
\textbf{Render} & \textbf{Tamaño} & \textbf{Estimado} \\
\midrule
Señal completa & 1\,800\,000 muestras & $\approx$336\,ms \estmark \\
Trial individual & 30\,501 muestras & $\approx$5.7\,ms \estmark \\
\midrule
\textbf{Total por clic en "Aplicar Filtro"} & (render) + (cómputo, \S6, medido: 44.9\,ms) & \textbf{$\approx$386\,ms} \estmark \\
\bottomrule
\end{tabular}
\end{center}

Si esta estimación es siquiera aproximadamente correcta, \texttt{filter} tiene el mismo problema que
\texttt{wavelet}: un render sin vectorizar que domina por completo el tiempo total de la operación (87\% del
estimado), muy por encima del costo real del filtro Butterworth en sí (13\%, medido). Es, junto con
\texttt{wavelet}/\texttt{wavelet\_\allowbreak average}, el otro candidato prioritario a vectorizar.

\subsection{Pipeline completo percibido por el usuario}

Sumando benchmarks medidos de distintas secciones (todos en el hilo principal, por lo que se perciben como
un solo bloqueo de la interfaz):

\begin{center}
\begin{tabular}{lr}
\toprule
\textbf{Flujo} & \textbf{Tiempo total} \\
\midrule
Cargar ABF + cortar trials & 78.4\,ms (medido) \\
+ Ver escalograma wavelet de 1 trial & $+$941\,ms (216\,ms CWT $+$ 725\,ms render, medidos) \\
\midrule
\textbf{Total: abrir archivo $\to$ ver 1 escalograma} & \textbf{$\approx$1.02\,s} \\
\bottomrule
\end{tabular}
\end{center}

% ==================================================================
\section{Resultados: \texttt{DataStore}}
% ==================================================================

\begin{center}
\begin{tabular}{lrl}
\toprule
\textbf{Benchmark} & \textbf{Mediana} & \textbf{Detalle} \\
\midrule
\texttt{add\_\allowbreak signal} ($\times$500) & 13.0\,ms & 0.026\,ms/señal; genera claves \texttt{raw\_\allowbreak signal\_\allowbreak N} con un \texttt{while} que revisa colisiones. \\
\texttt{set}$+$\texttt{get} ($\times$2000 pares) & 0.6\,ms & 0.00015\,ms/operación -- diccionario puro, O(1) como se espera. \\
\bottomrule
\end{tabular}
\end{center}

\texttt{DataStore} no es, ni de cerca, un cuello de botella incluso en escenarios con muchas más señales
cargadas de las que un usuario típico acumularia en una sesión. No requiere atención en 2.0.

\newpage
% ==================================================================
\section{Hallazgos priorizados}
% ==================================================================

\begin{enumerate}
  \item \critmark\ \textbf{Renderizado 2D sin vectorizar en \texttt{wavelet}/\texttt{wavelet\_\allowbreak
        average}} (725--742\,ms, hilo principal) -- el bloqueo de interfaz más grande y más fácil de
        eliminar de todo el informe: reemplazar el bucle
        \texttt{SetScalarComponentFromFloat} por \texttt{numpy\_\allowbreak to\_\allowbreak vtk} (mismo
        patrón que ya usa \texttt{core/utils/adapters.py}) puede reducirlo, según el micro-benchmark de
        \S\ref{sec:micro-bench}, a pocos milisegundos.
  \item \critmark\ \textbf{\texttt{wavelet.compute\_\allowbreak wavelet} en el hilo principal} (216\,ms) --
        el segundo bloqueo más grande. A diferencia del render, este es un problema algorítmico/de
        threading, no solo de vectorización: mover el CWT a un \texttt{QThread} (como ya hace
        \texttt{wavelet\_\allowbreak average}) resuelve el bloqueo de UI sin tocar la matemática.
  \item \critmark\ \textbf{Renderizado de línea sin vectorizar ni topar en \texttt{filter}} ($\approx$336\,ms
        estimado, hilo principal, sin tope de muestras) -- mismo patrón que (1), pero además sin el límite
        de seguridad (\texttt{MAX\_\allowbreak SAMPLES}) que sí tienen \texttt{erp}/\texttt{average}.
  \item \warnmark\ \textbf{\texttt{wavelet\_\allowbreak average.compute\_\allowbreak wavelet}}
        ($\approx$12.6\,s estimado para 60 trials) -- el número absoluto más grande del informe, pero de
        menor prioridad relativa porque ya corre fuera del hilo principal (QThread) y no congela la
        interfaz; sigue siendo candidato a paralelizar entre trials o cachear resultados parciales en 2.0.
  \item \warnmark\ \textbf{Inconsistencia de patrón en \texttt{average\_\allowbreak plugin.py}} (bucle sin
        vectorizar, pero topado a 2000 puntos, impacto $<$0.5\,ms) -- no es una prioridad de rendimiento, pero
        es la misma deuda de código que (1) y (3): en vez de 5 implementaciones distintas del mismo patrón,
        Gamma Lab 2.0 debería tener un único helper vectorizado para \emph{cualquier} conversión NumPy$\to$VTK
        (imagen 2D o tabla), reutilizado por todos los plugins.
  \item \okmark\ \textbf{Arranque del kernel, carga de ABF, corte de trials, \texttt{DataStore}, FFT/PSD} --
        todos por debajo de 80\,ms, ninguno requiere optimización en 2.0. No gastar presupuesto de ingeniería
        aquí.
\end{enumerate}

% ==================================================================
\section{Recomendaciones para Gamma Lab 2.0}
% ==================================================================

\begin{enumerate}
  \item \textbf{Un único helper vectorizado para imágenes 2D}, análogo a
        \texttt{core/utils/adapters.py} pero para \texttt{vtkImageData}: recibir una matriz \texttt{Z}
        (frecuencia $\times$ tiempo, o trial $\times$ tiempo) y construir el \texttt{vtkImageData} con
        \texttt{numpy\_\allowbreak to\_\allowbreak vtk(Z.ravel(order="C"))} $+$
        \texttt{img.GetPointData().SetScalars(...)}, en vez de que cada uno de \texttt{wavelet},
        \texttt{wavelet\_\allowbreak average} y \texttt{erp} reimplemente su propio doble \texttt{for}. Esto
        ataca los hallazgos (1) y (5) a la vez.
  \item \textbf{Usar el helper de tablas ya existente en \texttt{filter\_\allowbreak plugin.py}} en vez de su
        bucle propio de \texttt{InsertNextValue} -- ataca el hallazgo (3), y de paso agregar el mismo tope
        \texttt{MAX\_\allowbreak SAMPLES} que ya tienen \texttt{erp}/\texttt{average} como red de seguridad
        adicional.
  \item \textbf{Mover todo cómputo pesado a \texttt{QThread} de forma sistemática}, no solo donde ya se hizo
        (\texttt{wavelet\_\allowbreak average}): en particular \texttt{wavelet.compute\_\allowbreak wavelet}
        (216\,ms, hallazgo 2). El resto (FFT/PSD, $<$12\,ms) no lo necesita hoy, pero conviene decidir un
        criterio único (p.\,ej. ``todo cómputo $>$50\,ms va a QThread'') en vez de que sea caso por caso.
  \item \textbf{Conservar esta suite de benchmarks como guardia de regresión}, corriéndola en cada cambio
        significativo de rendimiento durante el desarrollo de 2.0 -- no solo al final. Ver
        sección~\ref{sec:comparacion-v2}.
  \item \textbf{Evaluar cachear o paralelizar el CWT entre trials} en \texttt{wavelet\_\allowbreak average}
        (hallazgo 4) -- a 209.5\,ms/trial, 60 trials en serie son $\approx$12.6\,s; con 16 núcleos lógicos
        disponibles en la máquina de referencia, paralelizar los CWT independientes entre trials (son
        embarrasingly parallel) podría reducir ese número de forma significativa sin cambiar el algoritmo.
  \item \textbf{Agregar un archivo \texttt{.edf} y/o \texttt{.mat} de referencia a \texttt{test/data/}} para
        poder benchmarkear (y validar contra MATLAB) los otros dos formatos que soporta
        \texttt{FileIOService}, hoy sin cobertura de ningún tipo.
\end{enumerate}

\newpage
% ==================================================================
\section{Metodología de comparación con Gamma Lab 2.0}
\label{sec:comparacion-v2}
% ==================================================================

Esta es la parte pensada explícitamente para dentro de un mes. El procedimiento:

\begin{enumerate}
  \item \textbf{No tocar \texttt{bench\_\allowbreak common.py} ni \texttt{compare\_\allowbreak report.py}} --
        son agnósticos a la versión de la app: solo cronometran funciones y comparan dos etiquetas del mismo
        JSON.
  \item \textbf{Copiar (o adaptar) los archivos \texttt{test\_\allowbreak *\_\allowbreak bench.py}} al
        repositorio de Gamma Lab 2.0. Si 2.0 mantiene los mismos nombres de módulo/función
        (\texttt{FileIOService.load\_\allowbreak abf}, \texttt{cut\_\allowbreak trials\_\allowbreak
        single\_\allowbreak channel}, \texttt{\_\allowbreak compute\_\allowbreak fft}, \texttt{compute\_\allowbreak
        wavelet}, \texttt{render\_\allowbreak scalogram}, \texttt{core.plugins.loader.discover}, etc.), los
        archivos deberían correr prácticamente sin cambios. Si 2.0 renombra o reorganiza módulos, el único
        ajuste necesario son los \texttt{import} de cada archivo -- la lógica de medición no cambia.
  \item \textbf{Correr contra 1.0 con una etiqueta clara} (ya hecho en este informe):
        \begin{verbatim}
$env:GAMMA_PERF_LABEL = "gammalab1_baseline_2026-08-16"
pytest test/performance_test -s
        \end{verbatim}
  \item \textbf{Correr la misma suite contra 2.0}, en la \textbf{misma máquina} si es posible (o dejando
        registrado el entorno de cada lado, capturado automáticamente por \texttt{record\_\allowbreak
        environment()} -- ver \S\ref{sec:entorno}):
        \begin{verbatim}
$env:GAMMA_PERF_LABEL = "gammalab2_baseline_2026-09-XX"
pytest test/performance_test -s
        \end{verbatim}
  \item \textbf{Comparar con el script ya existente:}
        \begin{verbatim}
python test/performance_test/compare_report.py `
    gammalab1_baseline_2026-08-16 gammalab2_baseline_2026-09-XX
        \end{verbatim}
        Esto imprime una tabla con el tiempo ``antes'', ``después'' y el factor de aceleración por cada
        benchmark que exista en ambas etiquetas, además del entorno de cada lado.
  \item \textbf{Repetir cada lado 2--3 veces} antes de sacar conclusiones sobre los benchmarks de
        renderizado VTK, por la variabilidad de primer render offscreen documentada en
        \S\ref{sec:resultados-render-lento}.
\end{enumerate}

\textbf{Qué preguntas debería poder responder esta comparación en 2.0:}
\begin{itemize}[nosep]
  \item ¿Se corrigió el patrón de renderizado sin vectorizar? (factor esperado: 65--150x en los benchmarks de
        \texttt{wavelet}/\texttt{wavelet\_\allowbreak average}/\texttt{filter}).
  \item ¿\texttt{wavelet.compute\_\allowbreak wavelet} sigue bloqueando el hilo principal, o se movió a
        background?
  \item ¿El arranque del kernel se mantuvo liviano, o una arquitectura nueva (más plugins, más servicios) lo
        volvió perceptible?
  \item ¿La carga de archivos y el corte de trials se mantienen submedio-segundo?
\end{itemize}

% ==================================================================
\section{Limitaciones conocidas}
\label{sec:limitaciones}
% ==================================================================

\begin{itemize}
  \item \textbf{\texttt{MainWindow} no se benchmarkeó.} Requiere una \texttt{QApplication} real con ciclo de
        eventos (\texttt{app.exec\_\allowbreak ()}), no solo un \texttt{vtkRenderWindow} offscreen como el
        resto de la suite -- automatizar esto de forma confiable (sin dejar procesos colgados) es trabajo
        adicional fuera del alcance de esta corrida. Es la única pieza de la secuencia de arranque real de
        \texttt{main.py} que falta medir.
  \item \textbf{\texttt{filter\_\allowbreak plugin.\_\allowbreak render\_\allowbreak filtered} no se
        benchmarkeó en vivo}, por la misma razón: construye un \texttt{QVTKRenderWindowInteractor} real y un
        \texttt{VTKContextMenu} que instancia un \texttt{QMenu} de Qt -- sin una \texttt{QApplication} viva,
        PyQt5 puede abortar el proceso en vez de lanzar una excepción Python capturable. La cifra de
        \S\ref{sec:estimaciones} es una \textbf{estimación} derivada del micro-benchmark sintético, no una
        medición directa.
  \item \textbf{\texttt{load\_\allowbreak edf}/\texttt{load\_\allowbreak mat} sin cobertura} -- no hay
        archivo de referencia en \texttt{test/data/}.
  \item \textbf{Los tiempos de renderizado VTK offscreen tienen variabilidad de proceso a proceso}, documentada
        en \S\ref{sec:resultados-render-lento} con un ejemplo real (5.70\,s vs.\ 725\,ms para el mismo
        código). Tratar una única corrida como definitiva, especialmente para los benchmarks de render, es
        un error metodológico -- de ahí la recomendación de repetir 2--3 veces en
        \S\ref{sec:comparacion-v2}.
  \item \textbf{Todas las cifras son de una sola máquina} (\S\ref{sec:entorno}). No se afirma que sean
        representativas de otro hardware -- son una línea base válida \emph{para comparar en esa misma
        máquina}, que es exactamente el uso previsto este mes.
\end{itemize}

% ==================================================================
\section{Tabla completa de resultados (apéndice)}
% ==================================================================

Todas las cifras de este informe, en un solo lugar, tal como quedaron guardadas en
\texttt{test/performance\_\allowbreak test/results/perf\_\allowbreak results.json} bajo la etiqueta
\texttt{gammalab1\_\allowbreak baseline\_\allowbreak 2026\allowbreak-08\allowbreak-16}.

\begin{center}
\begin{longtable}{L{9.5cm}rc}
\toprule
\textbf{Benchmark} & \textbf{Mediana} & \textbf{n} \\
\midrule
\endhead
\texttt{kernel.plugins.discover} (16 plugins) & 19.6\,ms & 5 \\
\texttt{kernel.plugin\_\allowbreak manager.load\_\allowbreak all} (16 plugins) & 20.2\,ms & 5 \\
\texttt{kernel.bootstrap\_\allowbreak completo} (16 plugins) & 19.9\,ms & 5 \\
\texttt{kernel.register\_\allowbreak plugin} (overhead, x200 ficticios) & 1.4\,ms & 10 \\
\texttt{fileio.load\_\allowbreak abf} (2x1\,800\,000) & 59.8\,ms & 5 \\
\texttt{trials.cut\_\allowbreak trials\_\allowbreak single\_\allowbreak channel} (30501x60) & 18.6\,ms & 10 \\
\texttt{filter.run\_\allowbreak filter} (señal completa, 1\,800\,000) & 43.2\,ms & 10 \\
\texttt{filter.run\_\allowbreak filter} (1 trial, 30501) & 1.8\,ms & 20 \\
\texttt{fft.\_\allowbreak compute\_\allowbreak fft} (todos los trials) & 3.8\,ms & 5 \\
\texttt{fft\_\allowbreak average.\_\allowbreak compute\_\allowbreak fft\_\allowbreak average} & 11.5\,ms & 5 \\
\texttt{psd.\_\allowbreak compute\_\allowbreak psd} (Welch) & 5.3\,ms & 5 \\
\texttt{psd\_\allowbreak average.\_\allowbreak compute\_\allowbreak psd} (Welch) & 9.2\,ms & 5 \\
\texttt{relative\_\allowbreak psd.\_\allowbreak compute\_\allowbreak psd} (Welch) & 6.3\,ms & 5 \\
\texttt{relative\_\allowbreak psd.\_\allowbreak compute\_\allowbreak relative\_\allowbreak psd} & 0.07\,ms & 5 \\
\texttt{wavelet.compute\_\allowbreak wavelet} (1 trial) & 215.6\,ms & 3 \\
\texttt{wavelet\_\allowbreak average.compute\_\allowbreak wavelet} (x20 trials) & 4190.2\,ms & 2 \\
\texttt{wavelet.render\_\allowbreak scalogram} (998x3051) & 725.3\,ms & 3 \\
\texttt{wavelet\_\allowbreak average.render\_\allowbreak scalogram} (998x3051) & 741.6\,ms & 3 \\
\texttt{erp.\_\allowbreak render\_\allowbreak heatmap} (60x30501$\to$60x1907) & 35.8\,ms & 3 \\
\texttt{adapters.dataset\_\allowbreak to\_\allowbreak vtk\_\allowbreak table} (1800000x3) & 21.0\,ms & 10 \\
\texttt{adapters.trials\_\allowbreak matrix\_\allowbreak to\_\allowbreak vtk\_\allowbreak table} (30501x61) & 11.1\,ms & 10 \\
\texttt{vtk\_\allowbreak array.bucle\_\allowbreak python} (50\,000 puntos) & 4.6\,ms & 5 \\
\texttt{vtk\_\allowbreak array.numpy\_\allowbreak to\_\allowbreak vtk} (50\,000 puntos) & 0.031\,ms & 5 \\
\texttt{vtk\_\allowbreak array.bucle\_\allowbreak python} (300\,000 puntos) & 28.4\,ms & 5 \\
\texttt{vtk\_\allowbreak array.numpy\_\allowbreak to\_\allowbreak vtk} (300\,000 puntos) & 0.432\,ms & 5 \\
\texttt{data\_\allowbreak store.add\_\allowbreak signal} (x500) & 13.0\,ms & 5 \\
\texttt{data\_\allowbreak store.set+get} (x2000 pares) & 0.6\,ms & 5 \\
\bottomrule
\end{longtable}
\end{center}

% ==================================================================
\section{Conclusión}
% ==================================================================

Gamma Lab 1.0 es, para la enorme mayoría de sus operaciones (arranque, carga de archivos, FFT, PSD,
\texttt{DataStore}), rápido: por debajo de 80\,ms, imperceptible para el usuario. El rendimiento de la
aplicación como un todo está, sin embargo, dominado por un número pequeño de puntos muy concentrados: el CWT
de wavelet (un problema de threading, ya resuelto parcialmente en \texttt{wavelet\_\allowbreak average} pero
no en \texttt{wavelet}) y el renderizado VTK sin vectorizar de los mapas 2D y de al menos un gráfico de línea
(\texttt{filter}) -- un problema que, a diferencia del anterior, \textbf{ya tiene una solución probada y en
uso dentro del mismo repositorio} (\texttt{core/utils/adapters.py}), solo que no se aplicó de forma
consistente a los cinco lugares que comparten el mismo patrón.

Esta suite de benchmarks y este informe quedan pensados para no ser un documento de una sola vez: la
sección~\ref{sec:comparacion-v2} deja el procedimiento exacto para que, durante el desarrollo de Gamma Lab
2.0 este mes, cada decisión de arquitectura pueda contrastarse contra esta línea base con
\texttt{compare\_\allowbreak report.py}, en vez de depender de percepciones subjetivas de velocidad.

\end{document}
